% Round 14. New extension of the R13 smooth germ lemma; independently cross-audited.
\section{Algebraic closed-immersion germs do not require smoothness}
\label{r14l:section}

The smooth analytic retraction used in Round 13 has a stronger
consequence for algebraic varieties. A singular closed subvariety
can be treated point by point through a smooth algebraic stratum.
The resulting statement is about the original actual sheaf on
the ordinary analytic topology, not about a replacement by forms
on a resolution.

Throughout this section $K$ has characteristic zero. Every variety
whose actual sheaf is written is reduced and pure-dimensional;
its analytification is good and boundaryless. Different varieties
in one morphism may have different dimensions. This purity
condition is retained because Mihatsch's Definition 4.2 is stated
for pure-dimensional spaces.

\begin{lemma}[Lifting a moment presentation across a closed immersion]
\label{r14l:local-lift}
For every algebraic locally closed immersion $j:T\hookrightarrow X$,
the actual restriction map
\[
                           j^{-1}\mathcal B_X\longrightarrow\mathcal B_T
\tag{L1}\label{r14l:restriction}
\]
is surjective on every stalk, in every tridegree.
No smoothness or normality assumption is needed.
\end{lemma}

\begin{proof}
Fix $x\in T^{\rm an}$ and an actual germ on $T$ represented
on an ordinary neighborhood $W$ by a finite unit tuple
$u=(u_1,\ldots,u_s)$ and an ambient Euclidean delta form $\gamma$.
Locally in $X$ the immersion is closed. Choose an ordinary
ambient neighborhood $O$ of $x$ such that
$O\cap T^{\rm an}\subset W$ and $T\cap O$ is closed in $O$.
Choose an affinoid neighborhood $V\subset O$ with $x$ in its
ordinary interior. The closed analytic subspace $T\cap V$ is
affinoid, with algebra $A/I$ if $V=\mathcal M(A)$.
The restrictions of the $u_i$ to $T\cap V$ lift to elements
$\widetilde u_i\in A$. Since
$|\widetilde u_i(x)|=|u_i(x)|>0$, their simultaneous nonvanishing
locus contains an ordinary neighborhood $V'$ of $x$ inside
$\operatorname{Int}(V)$. On $V'$ all $\widetilde u_i$ are
analytic units. The actual form
$t_{\widetilde u}^{\star}(\gamma)$ is a lift of the prescribed
germ by Mihatsch's genuine analytic pullback theorem. The
construction preserves the ambient tridegree.
\end{proof}

\begin{theorem}[Arbitrary algebraic immersion germ isomorphism]
\label{r14l:germ-iso}
For every algebraic locally closed immersion of the varieties
specified above, (L1) is an isomorphism of trigraded sheaves
with all actual differential operators. In particular,
\[
                          j^{-1}\mathcal B_X\simeq\mathcal B_T
\tag{L2}\label{r14l:germ-isomorphism}
\]
holds when either or both varieties are singular.
\end{theorem}

\begin{proof}
Only injectivity remains. Let an actual germ $\alpha$ on $X$
restrict to zero on $T$ at $x$. Locally embed the algebraic
variety $X$ as a closed subvariety of a smooth affine space $P$.
Lemma~\ref{r14l:local-lift} lifts the chosen germ of $\alpha$
to an actual germ $\widetilde\alpha$ on $P$.

Choose a smooth algebraic locally closed subvariety
$S\subset T$ whose analytification contains $x$. Such an $S$
always exists over a characteristic-zero field: take the reduced
Zariski closure in $T$ of the scheme point underlying $x$ and
then its nonempty smooth locus. That scheme point is the generic
point of its closure, hence belongs to this smooth locus.
Equivalently, use a finite algebraic smooth stratification of $T$.
After shrinking an algebraic ambient open containing $S$, the
immersion $S\hookrightarrow P$ is closed and both varieties
are smooth.

The restriction of $\widetilde\alpha$ to $S$ has zero germ at
$x$, since it factors through its zero restriction to $T$.
The Round 13 smooth closed-immersion germ lemma now implies
$\widetilde\alpha=0$ at $x$ in $\mathcal B_P$.
Restricting to $X$ gives $\alpha=0$. This proves injectivity.
All maps are genuine analytic pullbacks, so they preserve the
trigrading and commute with the actual differentials.
\end{proof}

\begin{corollary}[Ordinary degreewise extension and zero-neighborhood kernel]
\label{r14l:global-extension}
Let $j:T\hookrightarrow X$ be closed, with $X^{\rm an}$
paracompact and admitting ordinary smooth partitions of unity.
Then restriction on global actual forms is surjective in every
degree. Its kernel consists exactly of the forms which vanish
on an ordinary open neighborhood of $T^{\rm an}$.
\end{corollary}

\begin{proof}
Choose local moment lifts and multiply them by an ordinary
smooth partition of unity on $X$ subordinate to their open
neighborhoods and the complement of $T$. Their locally finite
sum, with the complement term chosen zero, lifts the prescribed
global form. This is degreewise extension only; multiplying by
a partition does not commute with $d'$ or $d''$.

The kernel assertion follows directly from (L2): zero
restriction means zero ambient germ at each point of $T$.
The union of the resulting zero neighborhoods is the required
ordinary neighborhood of $T^{\rm an}$.
\end{proof}

\subsection{Actual descent through a singular closed-locus square}

\begin{theorem}[Exact actual square for a proper modification]
\label{r14l:modification-square}
Let $q:Y\to X$ be a proper algebraic morphism which is an
isomorphism outside a reduced closed subvariety $Z\subset X$.
Put $E=(q^{-1}Z)_{\rm red}$, and write $i:E\hookrightarrow Y$,
$j:Z\hookrightarrow X$, and $q_E:E\to Z$. Assume that
$X,Y,Z,E$ are pure-dimensional varieties and that the ambient
analytic spaces admit ordinary smooth partitions of unity.
The following is degreewise exact as a sequence of actual
double complexes of global forms:
\[
0\longrightarrow\mathcal B(X)
 \xrightarrow{(q^*,j^*)}\mathcal B(Y)\oplus\mathcal B(Z)
 \xrightarrow{\ i^*-q_E^*\ }\mathcal B(E)
 \longrightarrow0.
\tag{L3}\label{r14l:modification-sequence}
\]
The center and exceptional variety may be singular; no normal
crossing or smooth-intersection hypothesis is used.
\end{theorem}

\begin{proof}
Surjectivity of the last arrow follows from
Corollary~\ref{r14l:global-extension} for $i$.
Suppose $(\beta,\gamma)$ lies in its kernel.
Choose a global form $\alpha_0$ on $X$ with
$j^*\alpha_0=\gamma$. Then
$\beta-q^*\alpha_0$ has zero restriction to $E$, so it
vanishes on an ordinary neighborhood $N$ of $E$.
Properness makes $q(Y\setminus N)$ closed and disjoint from
$Z$. Hence there is an ordinary neighborhood $V$ of $Z$ with
$q^{-1}(V)\subset N$.

Over $X\setminus Z$ transport $\beta-q^*\alpha_0$ through
the isomorphism $q$; it vanishes on $V\setminus Z$ and
therefore glues with the zero form on $V$. Call the resulting
global actual form $\alpha_1$. Then
$q^*\alpha_1=\beta-q^*\alpha_0$ and $j^*\alpha_1=0$.
Thus $\alpha_0+\alpha_1$ maps to $(\beta,\gamma)$.

If $(q^*\alpha,j^*\alpha)=0$, the first component makes
$\alpha$ zero outside $Z$ and the second makes it zero on an
ordinary neighborhood of $Z$. Hence $\alpha=0$.
Finally all arrows are actual pullbacks, so they commute with
both differentials. This proves (L3) as a short exact double
complex, not merely as a sequence of currents.
\end{proof}

\begin{corollary}[The same actual Aeppli obstruction for singular centers]
\label{r14l:aeppli}
Under the hypotheses of Theorem~\ref{r14l:modification-square},
there is an exact sequence
\[
\begin{split}
H_{\mathrm A}^{r-1,r-1}(Y)\oplus H_{\mathrm A}^{r-1,r-1}(Z)
 &\longrightarrow H_{\mathrm A}^{r-1,r-1}(E)
 \xrightarrow{\ \delta\ }H_{\mathrm{BC}}^{r,r}(X)\\
 &\longrightarrow H_{\mathrm{BC}}^{r,r}(Y)
                       \oplus H_{\mathrm{BC}}^{r,r}(Z).
\end{split}
\tag{L4}\label{r14l:aeppli-sequence}
\]
Every group here is the actual group on the original analytic
space. If $r>\dim Z$, the kernel of pullback to $Y$ is
identified with
\[
\frac{H_{\mathrm A}^{r-1,r-1}(E)}
 {i^*H_{\mathrm A}^{r-1,r-1}(Y)+
                         q_E^*H_{\mathrm A}^{r-1,r-1}(Z)}.
\tag{L5}\label{r14l:obstruction-quotient}
\]
\end{corollary}

\begin{proof}
Apply the direct lift-and-differentiate construction of the
Round 13 Aeppli connecting map to the short exact double
complex (L3). The proof uses only degreewise exactness and the
double-complex identities, so applies verbatim. In particular,
no purported general long exact sequence for Bott--Chern
cohomology is assumed. If $r>\dim Z$, actual forms on $Z$ in
bidegree $(r,r)$ vanish, so its Bott--Chern term is zero.
Exactness gives precisely (L5).
\end{proof}

\paragraph{Purity boundary and relation to the arrangement note.}
An arbitrary inverse image of a closed algebraic set need not
be pure-dimensional. We do not write $\mathcal B(E)$ for such
a space under Mihatsch's present definition. Formula (L3)
applies directly, for example, to a normalized blowup of a
pure-dimensional center when its reduced inverse image is an
effective Cartier divisor, and to modifications presented with
a pure Cartier closed locus. The compatible-family construction
for finite pure arrangements remains available when the chosen
union is not covered by this purity hypothesis. The new theorem
removes smoothness assumptions; it does not silently change the
definition on nonpure spaces.

\paragraph{Primary inputs.}
Mihatsch's Definition 4.2 and Theorem 4.5 supply the actual
local moment presentations and genuine pullback on
pure-dimensional boundaryless spaces:
\url{https://arxiv.org/html/2112.10018v4}.
The affinoid quotient property is the defining local property of
a closed analytic immersion. The smooth ordinary germ
isomorphism is the separately audited Round 13 theorem based on
Berkovich's degree-one etale germ criterion. The extension from
smooth to singular algebraic immersions is the new stratum
argument proved here.
